\section{Installation von Python und VSCodium}
Um mit dem Programmieren loslegen zu können, müssen wir zuerst die Programmiersprache auf unserem Computer installieren, sowie ein Programm, mit welchem wir Python-Code schreiben und ausführen können. Ein solches Programm wird typischerweise \gls{IDE} genannt. In diesem Skript verwenden wir die kostenfreie Programmiersprache \textit{Python}, sowie die ebenfalls kostenfreie, weit verbreitete \gls{IDE} \textit{VSCodium}. Der nachfolgende Abschnitt leitet Sie durch die Installation sowohl auf Windows wie auf MacOS.

\subsection{MacOS}
Um VSCodium unter MacOS zu installieren, benötigen Sie zuerst den Paket-Manager \texttt{brew}. Was macht \texttt{brew}? Laut der \href{https://brew.sh/de/}{offiziellen Webseite}: \enquote{Homebrew installiert Zeug, das du brauchst, das Apple aber nicht mitliefert.}

Falls Sie \texttt{brew} noch nicht auf Ihrem Mac installiert haben, tun Sie dies wie folgt:
\begin{enumerate}
	\item Öffnen Sie ein neues Terminal-Fenster, indem Sie zunächst die Spotlight-Suche mit \keys{\cmd+\space} (\keys{\cmd} +  Leertaste) öffnen in der Spotlight-Suche den Suchbegriff \texttt{Terminal} eingeben und die Suchanfrage mit \keys{\return} (ENTER-Taste) ausführen.
	\item Geben Sie folgende Code-Zeile ein und führen Sie diese aus (indem Sie mit der \keys{\return}-Taste bestätigen):
	      \begin{lstlisting}[language=bash, numbers=none]
/bin/bash -c "$(curl -fsSL https://raw.githubusercontent.com/Homebrew/install/HEAD/install.sh)"
\end{lstlisting}
\end{enumerate}

Führen Sie danach folgende Befehle aus, um VSCodium und Python zu installieren:
\begin{lstlisting}[language=bash, numbers=none]
brew install --cask vscodium
\end{lstlisting}
\begin{lstlisting}[language=bash, numbers=none]
brew install python3
\end{lstlisting}

Nun sollten sowohl VSCodium wie auch Python installiert sein. Falls Sie es nicht öffnen können und stattdessen eine Sicherheits-Warnmeldung erhalten, folgen Sie \href{https://support.apple.com/de-ch/guide/mac-help/mh40616/mac}{dieser Anleitung}, um das Programm dennoch zu öffnen.

\subsection{Windows}
Öffnen Sie das Programm \textit{Powershell} (Sie können mit dem \keys{\faWindows}-Knopf danach suchen) und geben Sie folgende Zeile ein, um den Package-Manager \texttt{winget} zu installieren:
\begin{lstlisting}[language=bash, numbers=none]
winget install -e --id Microsoft.PowerShell
\end{lstlisting}



Danach installieren Sie VSCodium und Python wie folgt:
\begin{lstlisting}[language=bash, numbers=none]
winget install -e --id VSCodium.VSCodium
\end{lstlisting}

\begin{lstlisting}[language=bash, numbers=none]
winget install -e --id Python.Python.3.11 --scope machine
\end{lstlisting}


Sie können nun mit dem \keys{\faWindows}-Knopf das Programm-Menu öffnen und nach \textit{VSCodium} suchen, welches nun installiert sein sollte.
